% Physical pages 3 and 4 of the concise study: the thematic orientation, in the
% order of the Mass, with what the Fathers say of each text. Every claim and
% quotation comes from the reviewed expansive study. The developed comparison
% of the three interpretations begins on the fifth page.

\runninghead{Themes and Movement}
\section{The Propers: Themes and Movement}
\zlabel{triptych:concise:themes:start}

In the Gospel of this Mass a householder will not let his servants pull up the
cockle that an enemy sowed among his wheat while men slept, and leaves both to
grow until the harvest. In the Introit God tells the exiles in Babylon that he
thinks thoughts of peace toward them and that they shall call upon him; in the
Epistle St Paul asks the Colossians to bear with and forgive one another and
to let the word of Christ dwell in them. What the Fathers say of these texts
gives three interpretations of the Mass. In the first, the separation of the
cockle from the wheat is left to the angels at the harvest, because those who
would weed may be deceived and the cockle may yet become wheat (St Augustine),
and Paul's patience and forgiveness are the servants' work meanwhile. In the
second, God, who could have done what he promised without the captives'
prayer, urges them to pray so that they may receive it (St Jerome), and the
people who sing the chants are those St Augustine calls captives of sin and
sinners crying from the deep. In the third, the enemy sows while men sleep, a
sleep that St John Chrysostom says can be avoided, and the watch against it is
kept by the word of Christ that Paul asks to dwell richly in the heart.

\subsection{Introit and Collect: the letter to the exiles, and a family that leans on grace}

The Introit comes from a letter that Jeremias sent from Jerusalem to the
exiles in Babylon (Jer 29:1--2). It tells them to build houses there, to
marry, and to ``seek the peace of the city, to which I have caused you to be
carried away captives; and pray to the Lord for it'' (v.~7). Then it promises
the return, when ``the seventy years shall begin to be accomplished in
Babylon'' (v.~10), and gives the words the antiphon sings in God's own voice.
The antiphon keeps three things God says he will do: think thoughts of peace,
hear those who call on him, and bring back the captivity. It leaves out that
he thinks peace ``to give you an end and patience,'' that he ``will be found''
by them, and the gathering ``out of all nations.'' God has promised the
return and still tells the captives to call on him. St Jerome writes that he
could have kept the promise without their prayer, but urges them to pray so
that they may be worthy to receive it (\work{In Hieremiam} V, PL~24
col.~893). The psalm verse gives thanks, ``Thou hast turned away the captivity
of Jacob,'' and St Augustine says it is spoken to Christ, who turned away
``our captivity, not by setting us free from the barbarians \dots\ but by
setting us free from bad works, from our sins'' (\work{Enarr.\ in Ps.}\ 84.2).

The Collect makes one petition, \latin{Familiam tuam \dots\ continua pietate
custodi}, ``keep thy family with continual mercy,'' for one purpose, that the
family may always be defended by God's protection. Between the two a relative
clause states the family's condition: \latin{quae in sola spe gratiae
caelestis innititur}, it leans on the hope of heavenly grace, and on that
alone. \latin{Innitor} is to lean one's weight on something.

\subsection{Epistle: what the elect of God put on}

Having told the Colossians to strip off ``the old man with his deeds'' (Col
3:9), St Paul names what the new man puts on: mercy, kindness, humility,
modesty, patience, ``bearing with one another and forgiving one another, if
any have a complaint against another. Even as the Lord hath forgiven you, so
do you also.'' Paul calls the offence a \latin{querela}, a complaint, and St
John Chrysostom says the word makes it ``a petty matter''; but by adding
``even as Christ forgave you'' Paul ``has persuaded us that even if we had
serious charges to bring, we ought to forgive'' (\work{Hom.\ in Col.}\ 8).
Above all these is charity, ``the bond of perfection.'' At verse~15 the Latin
reads \latin{pax Christi exsultet in cordibus vestris}, ``let the peace of
Christ rejoice in your hearts.'' The Greek that Chrysostom expounds has a verb
meaning ``let it act as umpire,'' and he builds on it a contest inside the
heart; St Thomas Aquinas, expounding the Latin, says the heart rejoices
because joy is the effect of charity and follows from peace (\work{Super
Col.}\ c.~3 lect.~3). Paul then asks that the word of Christ \latin{habitet in
vobis abundanter}, dwell in them abundantly, and that they teach one another
``in psalms, hymns and spiritual canticles, singing in grace in your hearts to
God.''

\subsection{Gradual and Alleluia: thanks taken from a lament, and a cry from the depths}

The Gradual sings two verses of Psalm~43: ``Thou hast delivered us, O Lord,
from them that afflict us: and hast put them to shame that hate us. In God we
will glory all the day.'' The psalm then turns to lament, ``But now thou hast
cast us off, and put us to shame'' (v.~10), and says ``for thy sake we are
killed all the day long'' (v.~22), which St Paul quotes of the Christians (Rom
8:36). Because of that quotation St Augustine hears in the psalm ``the voice
of the Martyrs,'' and he takes the past tense of the thanks as ``spoken of the
future under the figure of the past,'' because the future ``is as certain as if
it were past'' (\work{Enarr.\ in Ps.}\ 43.1, 6).

The Alleluia sings the opening of Psalm~129, the sixth of the penitential
psalms: ``From the depths I have cried to Thee, O Lord: Lord, hear my
prayer.'' The Missal has \latin{orationem meam}, ``my prayer,'' where the
Vulgate has ``my voice.'' St Augustine begins with Jonas, who cried from the
whale's belly, and then names the deep: ``this mortal life is our deep.'' The
one who cries from it is ``a sinner,'' and some are so far down that they ``do
not even cry from the deep'' (\work{Enarr.\ in Ps.}\ 129.1--2).

\subsection{Gospel: the cockle among the wheat}

Jesus tells the parable to the crowds from a boat, after the parable of the
sower (Mt 13:1--3). ``While men were asleep,'' an enemy sows cockle among the
good seed, and nothing shows until the blade has borne fruit. The cockle,
\latin{zizania}, is the darnel St Jerome calls \latin{lolium}, which in the
blade can hardly be told from wheat (\work{In Matthaeum} II, PL~26 col.~94).
The servants ask, ``Whence then hath it cockle?'' and offer to gather it, and
the householder refuses: ``No, lest perhaps gathering up the cockle, you root
up the wheat also together with it. Suffer both to grow until the harvest.''
Afterwards, in the house, Jesus tells the disciples that the field is the
world, the good seed the children of the kingdom, the enemy the devil, the
harvest the end of the world and the reapers the angels (vv.~37--39). That
explanation is not read in the Mass, and the Fathers read the parable with it.

Two details draw their comment. The first is the householder's reason. St
Augustine finds it in the servants, who may be deceived; Jerome in the
likeness of the two plants; St John Chrysostom in the bloodshed that weeding
would bring. Jerome and Chrysostom also read in it the hope that the cockle
may yet become wheat, and Augustine tells the tares the same. The second is
the sleep. The parable does not say who slept, and the Lord's explanation
passes over the sleepers. Jerome says they are ``the teachers of the
Churches'' (PL~26 col.~93); Chrysostom extends the danger to ``our rulers
\dots\ and not the rulers only, but the subjects too'' (\work{Hom.\ in Mt.}\
46.1).

\subsection{Offertory to Postcommunion: the cry again, the offering and the promise}

The Offertory sings the Alleluia's verses again and repeats the first line,
so that \latin{de profundis clamavi} is heard three times before the Secret.
The Secret names its gifts \latin{hostias \dots\ placationis}, offerings of
propitiation, and asks that God absolve our sins and \latin{nutantia corda tu
dirigas}, direct our wavering hearts: \latin{nutare} is to sway, as a thing
does that is not firmly set. Verse~4 of Psalm~129, which the Mass does not
sing, says ``with thee there is merciful forgiveness,'' Latin
\latin{propitiatio}, and Augustine asks, ``what is this propitiation, except
sacrifice?'' (129.3).

The Communion, Mk 11:24, was spoken the morning after Jesus cursed the fig
tree: whoever ``shall not stagger in his heart, but believe,'' can bid a
mountain be cast into the sea, and ``therefore'' whatever they ask in prayer,
believing, they shall receive. The next verse, which the antiphon does not
sing, says: ``when you shall stand to pray, forgive, if you have aught against
any man'' (v.~25). The Postcommunion calls what has been received
\latin{per haec mysteria} a \latin{pignus}, a pledge, and asks for the
\latin{effectus}, the effect, of the salvation it pledges.
\zlabel{triptych:concise:themes:end}
